\section{Conclusion}

In this paper, we address the LLM ``refusal dilemma'' through solvability boundary learning, separating input-side answerability defects from model-side inability on well-posed tasks. Using our UnsolvableQA benchmark and UnsolvableRL framework, we demonstrate that recognizing missing-information and constraint-conflict defects acts as a powerful {reasoning regularizer}. This paradigm enables robust unsolvability detection ($U > 85\%$), sharper boundaries on solvable tasks, two-family defect diagnosis, and zero-shot clarification transfer on IN3; 4B/8B models further achieve OOD gains on GPQA (+12.0\%) and ReliableMath (+12.1\%). By turning refusal penalties from capability collapse into a rigor catalyst, our work distinguishes intrinsic query flaws from model limitations.

\section*{Limitations}
Despite its diversity, our study focuses on structured missing-information and constraint-conflict defects, and does not cover the full spectrum of subjective ambiguity, undefined user intents, or subtle rule-level violations. Although diagnostic tuning shows zero-shot transfer to clarification behavior on IN3, our training objective remains primarily detection-oriented. A full closed-loop repair agent that optimizes when to ask, how to integrate user responses, and whether the repaired task is successfully solved remains future work.% Finally, while validated up to the 8B scale, computational constraints prevented testing on massive architectures (e.g., 70B+), which remains a subject for future verification.
